\documentclass[fleqn,10pt]{wlscirep}
\usepackage[T1]{fontenc}
\usepackage[utf8]{inputenc}
\usepackage{mathptmx}
\usepackage{booktabs,dcolumn}
\usepackage{color, colortbl}
\usepackage{graphicx}
\usepackage{hyperref}

\begin{document}

\thispagestyle{empty}

\begin{center}
{\LARGE \textbf{Cover Letter}}
\end{center}

\vspace{1cm}

\textbf{Date:} \today

\textbf{To:} The Editors \\
\textit{Scientific Reports}

\vspace{0.5cm}

\textbf{Dear Editors,}

We are pleased to submit our manuscript entitled \textbf{``Discovery of Diagnostic and Prognostic Long Non-Coding RNA Biomarkers in Patients with Colorectal Cancer Using Machine Learning Algorithms''} for consideration for publication in \textit{Scientific Reports}.

\subsection*{Summary of the Work}

This study develops an integrated machine learning pipeline for lncRNA biomarker discovery in colorectal cancer. Using TOIL-recomputed TCGA-GTEx RNA-seq data (n = 639) and TCGA-COAD clinical follow-up (n = 443), we:

\begin{enumerate}
    \item Identified 36 lncRNAs achieving near-perfect tumor-normal discrimination (SVM AUC = 0.999, random forest AUC = 0.997)
    \item Developed a 97-feature Cox PH prognostic model that improved Harrell's C-index by +0.123 over stage alone
    \item Externally validated the prognostic risk score across three independent cohorts totaling 917 patients (TCGA-READ, GSE39582, GSE17536), demonstrating consistent hazard ratio directionality
    \item Demonstrated that pathological stage remains the dominant prognostic factor (RF VIMP = 1.000), with lncRNAs providing incremental value
\end{enumerate}

\subsection*{Key Strengths}

\begin{itemize}
    \item Five classification algorithms benchmarked with rigorous cross-validation
    \item External validation across RNA-seq and microarray platforms spanning 917 patients
    \item Comprehensive cross-cohort clinical comparison with statistical testing
    \item Transparent discussion of platform compatibility limitations and annotation drift
    \item All analysis code publicly available at \url{https://github.com/woodhaha/CRCproject-lncrna}
\end{itemize}

\subsection*{Significance}

This work addresses a critical gap in lncRNA biomarker research: the systematic integration of diagnostic classification and prognostic modeling with multi-cohort external validation. Our findings highlight important methodological considerations—GENCODE annotation stability, cross-platform gene coverage, and the distributed nature of lncRNA prognostic signals—that should inform future biomarker discovery studies.

\subsection*{Declarations}

\begin{itemize}
    \item This manuscript has not been published and is not under consideration elsewhere
    \item All authors have read and approved the manuscript
    \item The authors declare no competing interests
    \item All data sources are publicly available as described in Methods
    \item Analysis code is available at \url{https://github.com/woodhaha/CRCproject-lncrna}
\end{itemize}

\vspace{0.5cm}

Thank you for your time and consideration. We look forward to your response.

\vspace{0.5cm}

\textbf{Sincerely,}

\vspace{0.3cm}

Guanyu Wang, MD \\
Department of General Surgery \\
Sir Run Run Shaw Hospital, School of Medicine \\
Zhejiang University \\
Hangzhou, Zhejiang 310016, China \\
Email: wangguanyu@zju.edu.cn

\end{document}
